\documentclass[12pt,fleqn]{amsart}

\usepackage[T1]{fontenc}
\usepackage[utf8]{inputenc}
\usepackage{graphicx}
\usepackage{float}
\usepackage{xcolor}
\usepackage[colorlinks = true, linkcolor = blue]{hyperref}
\usepackage{url}

%Reset the mathcal
\DeclareMathAlphabet{\mathcal}{OMS}{cmsy}{m}{n}
\usepackage{amsmath,amssymb,amsthm,textcomp}
\usepackage{multicol}
\usepackage{tikz}
\usepackage[normalem]{ulem}
\usepackage{bbm} % using for the characteristic function

\usepackage{geometry}
\geometry{left=15mm,right=15mm,%
bindingoffset=0mm, top=20mm,bottom=20mm}
\linespread{1.3}
\newcommand{\linia}{\rule{\linewidth}{0.5pt}}

\usepackage[open,openlevel=1]{bookmark}

%\newtheorem{theorem}{Theorem}
\newtheorem{theorem}{Theorem}[section]
\renewcommand{\thetheorem}{\arabic{section}.\arabic{theorem}}
\newtheorem{corollary}{Corollary}[theorem]
\newtheorem{lemma}[theorem]{Lemma}
\newtheorem{proposition}[theorem]{Proposition}
\theoremstyle{definition}
\newtheorem{definition}{Definition}[section]
\newtheorem{assumption}{Assumption}[section]
\renewcommand{\thedefinition}{\arabic{section}.\arabic{definition}}
\newtheorem*{example}{Example}
\newtheorem{remark}{Remark}
\allowdisplaybreaks

%algorithm
\usepackage{algorithm}
\usepackage{algorithmic}


%% SYMBOLS
%divergence
\AtBeginDocument{
  \let\div\relax
  \DeclareMathOperator{\div}{div}
}
%contradiction
\usepackage{stmaryrd}
\usepackage{marvosym}
%range
\AtBeginDocument{
  \let\ran\relax
  \DeclareMathOperator{\ran}{ran}
}
%ess span dist
\DeclareMathOperator*{\esssup}{ess\, sup}
\DeclareMathOperator*{\spa}{span}
\DeclareMathOperator*{\dist}{dist}
\DeclareMathOperator*{\atom}{atom}
\DeclareMathOperator*{\cont}{cont}
\DeclareMathOperator*{\sign}{sign}
\DeclareMathOperator*{\supp}{supp}
\DeclareMathOperator*{\dom}{dom}
\DeclareMathOperator*{\cl}{cl}
\DeclareMathOperator*{\relu}{ReLU}


\title[Review]{Review}
\begin{document}
\section{Abstract}
In many nonlinear dynamic control problems, the optimal value function is the solution of 
the Hamilton-Jacobi-Bellman PDEs, which are often high-dimensional. 
The data-driven approach with gradient augmentation has been used to solve the high-dimensional PDEs in order to circumvent the curse of dimensionality. 
In \cite{azmi2021optimal} (Kunisch et al., 2021), a data set consisting of state-adjoint-control triples is generated by solving the open-loop problem provided by the Pontryagin Maximum Principle. 
The optimal value function is in turn approximated by hyperbolic cross basis functions with $l^1$ penalty to enhance the sparsity.
The gradient augmentation combined with hyperbolic cross basis is proven to be very effective to overcome the curse of dimensionality. \par 
 
In terms of different approximation methods besides the hyperbolic basis functions, there is a growing literature on solving the nonlinear optimal control problem using neural networks. 
Although the number of neurons may not serve as a control parameter of the model capacity \cite{neyshabur2014search}, sparsity is crucial for the computational efficiency in many problems. 
Furthermore, it is also necessary for certain applications, such as deconvolution of point-like objects. 
 A sparse learning procedure usually adopts convex reformulation of the problem and employ a sparse promoting convex penalty term on the outer weights, a popular choice is the $l^1$ function (LASSO). 
 However it was indicated in \cite{pieper2022nonconvex} that the convex penalty terms are not able to effectively eliminate the redundancy, especially when the neurons are clustering in a small neighborhood. 
 A nonconvex sparse learning method is proposed by the author. \par
In the paper \cite{pieper2022nonconvex}(Pieper et al., 2022), the authors choose ReLU as the activation function. 
The ReLU functions enjoy many nice properties such as Lipschitz continuity and positive homogeneity, which are essential to prove the optimality and for the penalty term to promote the sparsity. 
The problem is that the ReLU activation functions are not smooth, which may suffer from exploding gradient. 
In this thesis we propose different activation functions that may circumvent this problem. 
Furthermore, we investigate a different nonconvex penalty that further enhaces the sparcity. 
And as in \cite{azmi2021optimal} we incorporate the gradient information in the cost function to obtain optimal convergence.\par  

 
\section{Introduction}
The problem under consideration is the dynamic optimal control problem in the space-time domain
\begin{equation}
\label{opt}
\tag{$P$}
\begin{aligned}
	&\inf_{u \in L^2(t_0, T, \mathbb{R}^k)} J(u; t_0, y):= \int_{t_0}^T l(y(t)) + \beta |u(t)|^2 dt, \quad \beta >0\\
	&\text{s.t.}\\
	& \dot{y} = f(y) + g(y) u,\\
	& y(t_0) = x,
\end{aligned}
\end{equation}
where we assume $y(t) \in \mathbb{R}^d, u(t) \in \mathbb{R}^k$, $l \in C^1(\mathbb{R}^d, \mathbb{R}), f \in C^1(\mathbb{R}^d, \mathbb{R}^d), g \in C^1(\mathbb{R}^d, \mathbb{R}^{d \times k})$. We assume the minimizer exists.\par

The first method to solve the optimization problem involves computing a state-adjoint-control triple originating from the initial condition $x$, known as the open-loop control. The Pontryagin's Maximum Principle(PMP) yields the first-order optimality condition in the form of a forward-backward coupling Hamilton Jacobi ODE:
\begin{equation}
\label{TPBVP}
\tag{TPBVP}
\left\{
\begin{aligned}
	&\frac{d}{dt}y^* = f(y^*) + g(y^*)u^*,\\
	& y^*(t_0) = x, \\
	& -\frac{d}{dt} p^* = \partial_y \bigg( f(y^*) +  g(y^*) u^* \bigg)^Tp^* + \partial_y l(y^*),\\
	& p^*(T) = 0.
\end{aligned}
\right.
\end{equation}
By the maximality condition of the Hamiltonian at the optimal control we get the optimality condition
\begin{equation}
	u^* = - \frac{1}{2\beta} g(y^*)^T p^* \quad \forall t \in (t_0, T). 
\end{equation}

The optimal trajectories obtained from the solution of the open-loop system are not robust with respect to disturbance \cite{azmi2021optimal}, since the existing numerical methods for two-point boundary value problem in general tend to diverge. The closed-loop method, on the other hand, is based on the dynamic programming. By introducing the value function 
\begin{align*}
	 V(t, y):= \inf_u J(u; t, y) \quad \text{ subject to the dynamic system},
\end{align*}
we synthesise the control by the maximum principle:
\begin{align*}
	u^*(t, x) = -\frac{1}{2\beta} g(y^*)^T p^* ,
\end{align*} 
and get a nonlinear hyperbolic partial differential equation, the Hamilton-Jacobi-Bellman (HJB) equation:
\begin{equation}
\label{HJB}
\tag{HJB}
\left\{
\begin{aligned}
	& \partial_t V - \frac{1}{4\beta} \nabla V(t, x)^T  g(x) g^T(x) \nabla V (t, x) + \nabla V(t, x)^T f(x) + l(x)  = 0 \\
	& V(T, x) = 0. 
\end{aligned}
\right.
\end{equation}
In this case the optimal control problem becomes solving a nonlinear hyperbolic partial differential equation. 
The method is more robust compared to the open-loop method but suffers from 2 main drawbacks. 
First, in the case where the dimension of the state space is large, we encounter the \textit{curse of dimensionality}(\textit{Bellman)} and the computational cost is formidable if it is not impossible. 
Second, since we usually linearise the system locally and solve the Riccati equation, the optimal control differs from the Riccati-based approximation especially in the transient phase. \par  

Many existing methods are to overcome the curse of dimensionality, in particular, the data-driven approaches. 
A data-driven approach usually approximates the optimal value function on a causality-free data set of state-value pair. 
The model does not require a grid and can be applied to high dimensional problem. 
In particular, there is a growing literature on solving the high-dimensional problem using  neural network. 


\subsection{theoretical review and notations}
Throughout the paper we denote $\Omega$ as $\mathbb{R}^{d + 1}$. 
A finite Radon measure $\mu$ is a $\sigma$-additive $\mathbb{R}^m$-valued set function defined on the Borel sets of $\Omega$ with $\mu(\emptyset) = 0$.                                                                              
The \emph{total-variation measure} of a finite Radon measure $\mu$ is the positive measure defined by                                                                                                                                       
\[                                                                                                                                                                                                                                          
|\mu|(E) = \sup \left\{ \sum_{n=1}^{\infty} |\mu(E_n)| \;:\; \bigcup_{n=1}^{\infty} E_n = E,\; \{E_n\} \text{ disjoint and measurable} \right\}.
\]
We denote the space of finite Radon measures by $M(\Omega, \mathbb{R}^m)$ (and $M(\Omega)$ if $m = 1$). It can be normed by the total variation norm
\[
\|\mu\|_{M} = |\mu|(\Omega).
\]
Each finite Radon measure $\mu$ possesses a \emph{polar decomposition}, \emph{i.e.}, there exists a function $\sigma \in L^1_{|\mu|}(\Omega, \mathbb{R}^m)$ with $|\sigma(x)| = 1$ for $|\mu|$-almost all $x \in \Omega$ such that
\[
\mu(E) = \int_E \sigma \, \mathrm{d}|\mu| \quad \text{for all measurable } E \subset \Omega.
\]
An important result about the space $M(\Omega, \mathbb{R}^m)$, which is known as the Riesz theorem, is its characterization as a dual space. To state this theorem, introduce the space
\[
C(\Omega, \mathbb{R}^m) = \overline{\{u \in C(\Omega, \mathbb{R}^m) \;:\; \operatorname{supp} u \Subset \Omega\}}^{\|\cdot\|_\infty}, \quad \|u\|_\infty = \sup_{x \in \Omega} |u(x)|_*
\]
which is the completion of continuous compactly supported functions with respect to the supremum norm. 
For $m = 1$ we denote it by $C(\Omega)$. 
The Riesz theorem now states that the dual space $(C(\Omega,
\mathbb{R}^m))^* = M(\Omega, \mathbb{R}^m)$, which also implies the completeness of the space of finite Radon measures. The corresponding dual pairing is defined by
\[
\langle \mu,\, u \rangle = \sum_{i=1}^{m} \int_\Omega u_i \, \mathrm{d}\mu_i
\]
for each $\mu \in M(\Omega, \mathbb{R}^m)$ and $u \in C(\Omega, \mathbb{R}^m)$.
Without further specification, we denote $\langle \cdot, \cdot\rangle$ the dual bracket $\langle \cdot, \cdot \rangle_{M(\Omega), C(\Omega)}$. 

We also introduce the following notation:
\begin{itemize}
	\item $M^0(\Omega)$: the space of all nonatomic measures on $(\Omega)$. 
	\item $M^\sharp(\Omega)$: the space of all purely atomic measures on $\Omega$. 
\end{itemize}

We decompose $\mu \in M(\Omega)$ by $\mu = \mu_{\cont} + \mu_{\atom}$, where $\mu_{\cont} \in M^0(\Omega)$ has support $\cont \mu$ and $\mu_{\atom} \in M^\sharp(\Omega)$ has support $\atom \mu$. Therefore we have
\begin{align*}
	M(\Omega) = M^0(\Omega) \oplus M^\sharp(\Omega). 
\end{align*}



\section{the sparse learning problem}
In the setting of supervised learning we use the neural network to approximate the optimal value function:

\begin{equation}
	\mathcal{N}_{\omega, c}(x) := \sum_{n = 1}^N c_n \sigma (a_n \cdot x + b_n) \in \mathbb{R},
\end{equation}

where $\sigma$ is a sufficiently smooth activation function and $N$ the width of the network. And the gradient is

\begin{equation*}
	\nabla_x  \mathcal{N}_{\omega, c}(x) =  \sum_{ n = 1}^N c_n \nabla_x \sigma(a_n \cdot x + b_n) \in \mathbb{R}^d .
\end{equation*}

In \cite{azmi2021optimal} the author explored the relationship between (\ref{TPBVP}) and (\ref{HJB}) and suggested a gradient-augmented approach. 
In the data generating process, the adjoint state $p^*$ is obtained as a byproduct of solving (\ref{TPBVP}). 
Moreover, the adjoint state provides the gradient information of the optimal value function: $p^*(t) = \nabla V(t, y^*(t))$ (\cite{cristiani2010initialization} \cite{cannarsa1991some} \cite{clarke1987relationship}). 
Instead of state-value pair we get the dataset of the triple $\{x^i, V(x^i), \nabla V(x^i)\}_{i = 1}^K$ that is enriched with gradient information. 
We have the minimization problem
\begin{equation}
\label{lost_dis}
\begin{aligned}
	&\min_{N \in \mathbb{N}, (a_n, b_n, c_n) \in \mathbb{R}^d \times \mathbb{R} \times \mathbb{R}, } l_K (\mathcal{N}_{\omega, c}; V),\\
	& \text{with}\\
	& l_K (\mathcal{N}_{\omega, c}; V) := \frac{1}{2K} \bigg( \sum_{k = 1}^K  |V(x^k) - \mathcal{N}_{\omega, c} (x^k)|^2 +  \sum_{k = 1}^K \sum_{i = 1}^d | \partial_{x_i} V(x^k) - \partial_{x_i} \mathcal{N}_{\omega, c} (x^k) |^2 \bigg),
\end{aligned}
\end{equation}
where $x^k$ is a sample of the random variable $x \in L^2(D, \nu)$ with probability distribution $\nu$ on the domain $D$. 
$K$ can be potentially large. The loss function can be understood as the empirical estimation of the $H^1$ loss
\begin{align*}
	l_\nu (\mathcal{N}_{\omega, c}; V) = \frac{1}{2} \int_D |V(x) - \mathcal{N}_{\omega, c} (x)|^2 + \sum_{i = 1}^d | \partial_{x_i} V(x) - \partial_{x_i} \mathcal{N}_{\omega, c} (x) |^2 d\nu .
\end{align*}



\subsection{the activation functions}
First we discuss the choice of activation functions. 
The activations function we consider satisfy the following assumption.
\begin{assumption}[on activation function]
\label{assumption on actfun}
\quad
\begin{enumerate}
\item  $\sigma$ is differentiable w.r.t. $x$. 
\item There exists $ \Lambda > 0$, such that for all $x \in D$, $\omega_1, \omega_2 \in \Omega$, it holds 
\begin{align*}
	&|\sigma(x, \omega_1) - \sigma(x, \omega_2) | \le \Lambda (1 + \|x\|) |\omega_1 - \omega_2|,\\
	& |\partial_{x_i} \sigma(x, \omega_1) - \partial_{x_i} \sigma(x, \omega_2) | \le \Lambda (1 + \|x\|) |\omega_1 - \omega_2| \text{ for all $i = 1, \dots, d$. } 
\end{align*}
\end{enumerate}
\end{assumption}

The Lipschitz continuity of each partial derivative is also required since we are including the gradient information of the value function.



\subsubsection*{The $p$-ReLU function}
\quad\par
First we consider the ReLU function with the power $p$:
\begin{align*}
	\sigma(x, \omega) = \max\{0, \omega \cdot x\}^{p}, \text{ with } p \in (2, \infty). 
\end{align*}
In this case the gradient is 
\begin{align*}
	\nabla_x \sigma(x, \omega) = 
	\begin{bmatrix}
		p \max\{0, \omega \cdot x \}^{p - 1} \omega_1\\
		p \max\{0, \omega \cdot x \}^{p - 1} \omega_2 \\
		\dots \\
		p \max\{0, \omega \cdot x \}^{p - 1} \omega_d
	\end{bmatrix} \in \mathbb{R}^d. 
\end{align*}

The activation function is non-decreasing and positive homogeneous of degree $p$, and it admits the following nice properties.
\begin{itemize}
	\item Since they are not polynomial, they can linearly approximate any measurable function \cite{leshno1993multilayer}.
	\item They are invariant w.r.t. the change of scale of data, which we will explain further in detail. 
\end{itemize}

For $p > 1$,  the activation function not only preserves the nice property such as homogeneity of ReLU, but also smooths the kinks. We gain 2 things from the differentiability: 1. It circumvents the non-differentiable corners of the ReLU function; 2. The gradient information of the activation function can be used to calculate the gradient-augmented loss function. More generally, we assume the activation function to satisfy the following condition:

\begin{assumption}[p-homogeneity]
	$\sigma$ is $p$-homogeneous w.r.t. $\omega$ for some $p > 2$, i.e., for all $x \in D$, $\lambda > 0$ it holds
\begin{align*}
	\sigma(x, \lambda \omega) = \lambda^p \sigma (x, \omega)\quad \text{ for all } \omega \in \mathbb{R}^{d + 1}.
\end{align*}
\end{assumption}
\begin{remark}
For $p > 2$, the activation function is twice differentiable and the gradient-augmented optimization problem is a smooth optimization problem.
The value part of the loss is approximated by a high-order ReLU activation function $\text{ReLU}^{2 + \epsilon}$, and the gradient by $\text{ReLU}^{1 + \epsilon}$,  for a small $\epsilon$. In the case of $p = 2$, the optimization problem is still well defined but we would have non-differentiable corners for $\partial_{x_i} \sigma$ as in the ReLU case. For $p < 2$, the Lipschitz continuity of the gradient function is compromised. 
\end{remark}



\subsubsection*{The $tanh$ activation function}



\subsection{the nonconvex penalty}
As introduced in the abstract, the $l^1$ penalty fails to detect the redundancy of the neurons clustering in a small neighborhood, which motivates the nonconvex penalty term. First we introduce the following assumption:
\begin{assumption} \label{pen1}
	The penalty term $\phi: \mathbb{R}_+ \rightarrow \mathbb{R}_+$ satisfies: $\phi$ is differentiable at 0 with $\phi'(0) = 1$. 
\end{assumption}

\begin{assumption} \label{pen2}
The penalty term $\phi: \mathbb{R}_+ \rightarrow \mathbb{R}_+$ is differentiable, concave, nondecreasing and coercive($\phi(z) \rightarrow \infty$ as $z \rightarrow \infty$) with $\phi(0) = 0$.
\end{assumption}
 
The Assumption \ref{pen2} describe the global property of the penalty functions. 
By introducing the nonconvexity the coercivity of the quadratic cost is compromised. 
With Assumption \ref{pen1}, \ref{pen2} the optimization problem stays coercive such that we can apply the direct method of calculus of variation. 
The following lemma follows from the assumption. 
 
\begin{lemma} If Assumption \ref{pen2} is satisfied, then it holds
\begin{itemize}
	\item subadditivity: for all $a, b$,  $\phi(a + b) \le \phi(a) + \phi(b)$. 
\end{itemize}
Furthermore, if Assumption \ref{pen1} is satisfied, it holds
\begin{itemize}
	\item $ \phi(|z|) \le |z|$.
\end{itemize}
\end{lemma}
\begin{proof}
	\quad\\
	\textbf{1.}  It holds
	\begin{align*}
		\phi(z_1 + z_2) &= \int_0^{z_1} \phi'(z) dz + \int_0^{z_2} \phi'(z_1 + z) dz, \\
(\text{concavity)}\quad		&\le  \int_0^{z_1} \phi'(z) dz + \int_0^{z_2} \phi'(z) dz, \\
\quad		& = \phi(z_1) + \phi(z_2).
	\end{align*}
	\quad\\
	\textbf{2.} it holds for $z \ge 0$:
	\begin{align*}
		\phi(z_1) &= \int_0^{z_1} \phi'(z) dz\\
		 (\phi \text{ concave and } \phi'(0) = 1)  & \le \int_0^{z_1} 1 dz \\
		 &= z_1
	\end{align*}
\end{proof}
 
 
Assumption \ref{pen1}, \ref{pen2} ensures that the nonconvex problem to admits a global minimizer. 
To enhance the sparsity, we need to exert strong penalizing effect on the weight on 2 neighboring neurons such that merging them would in fact reduce the cost. 
We are mainly interested in the behavior of $\phi'$ in a neighborhood of $0$: 
 
\begin{assumption}[Local concavity]
 \label{pen3}
 The penalty term $\phi: \mathbb{R}_+ \rightarrow \mathbb{R}$ satisfies:
 \begin{enumerate}
  	\item There exists $\gamma_1 \ge 0, \hat{z}_1 > 0$, such that $  -\gamma_1(z_2 - z_1) \le \phi'(z_2) - \phi'(z_1) \le 0$ for all $0\le z_1 \le z_2 \le \hat{z}_1. $ ($\gamma$-convex)
 	\item There exists $\gamma_2 > 0$, $\hat{z}_2 > 0$, such that $\phi'(z_2) - \phi'(z_1) \le -\gamma_2(z_2 - z_1) \le 0 $ for all  $0 \le z_1 \le z_2 \le \hat{z}_2. $ (strongly concave)
 \end{enumerate}
 \end{assumption}
 \begin{remark}
 The Assumption \ref{pen3} can be interpreted as the concave function being strongly concave but exhibits also certain convexity locally at 0. It is worth noticing that it is entirely a local property that only needs to take effect at 0. Assumption \ref{pen3}(2) induces  strict subadditivity that takes effect at small weights. \par
 In general the global concavity in the Assumption \ref{pen2} doesn't imply locally uniform concavity in the Assumption \ref{pen3}(2). 
 Observe, for example $\phi(x) =  x - \frac{2}{3}x^{\frac{3}{2}}$, where the function satisfies the Assumption \ref{pen2} but the locally uniform concavity fails. However since Assumption \ref{pen3} is only local property we may construct piecewise penalty by first constructing the penalty around the neighborhood of 0. 
  
 \end{remark}

\begin{proposition}[the second order estimate]
\label{secondorderestimate}
	Given $\phi$ that satisfies the Assumption \ref{pen1}, \ref{pen2}, if $\phi$ satisfies further Assumption \ref{pen3}, then it holds
	\begin{enumerate}
		\item for all $ z \in [0, \hat{z}_1]: \phi(z) \ge z - \frac{\gamma_1}{2} z^2$, 
		\item for all $z \in [0, \hat{z}_2]: \phi(z) \le z - \frac{\gamma_2}{2}z^2$.
	\end{enumerate} 
\end{proposition}

By the proposition above we know any penalty satisfies Assumption \ref{pen1} and \ref{pen3} admits such upper and lower bound. 
The function
\begin{equation}
	\phi_{\log, \gamma}(z) = \frac{1}{\gamma} \log ( 1 + \gamma z)
\end{equation}
for $\gamma > 0$ and its convex combination with the $l1$ norm will be the main function of choice for our numerical analysis. 



\subsection{the integral representation}
We use the shallow neural network to approximate the value function because its theoretical property is well-studied. 
The nice properties of the integral representation of the neural network has been extensively investigated. In \cite{bengio2005convex} and \cite{neyshabur2014search} the convex nature of the shallow neural network has been proved. 
In \cite{ongie2019function} and \cite{savarese2019infinite} the representability of the shallow ReLU neural network is defined, the functional class, which can be represented by the integral neural network with finite Radon measure, is specified.  
We introduce the following notation: $\Omega$ a subset of $\mathbb{R}^{n + 1}$, $D \subset \mathbb{R}^d$, $M(\Omega)$ the measure on $\Omega$ with bounded variation. For $\mu \in M(\Omega)$, we define the operator $\mathcal{N}: M(\Omega) \rightarrow L^2(D, \nu) $: 
\begin{align*}
	[\mathcal{N}\mu] (x) = \int_\Omega \sigma (\omega , x) d\mu(\omega).
\end{align*}

Since for all $x$, $\sigma(\cdot, x)$ is bounded by a $\mu$-integrable function , and for all $\omega$, $\sigma(\omega, \cdot )$ is continuous, we can differentiate under the integral sign and get
\begin{align*}
	\nabla_x [\mathcal{N}\mu](x) = \int_\Omega \nabla_x \sigma(\omega,  x) d\mu(\omega). 
\end{align*}

To abuse the notation, we define the operator
\begin{align*}
	\nabla \mathcal{N}: M(\Omega) \rightarrow L^2(D, \nu)^d, \quad [\nabla\mathcal{N} \mu](x) = \nabla_x[\mathcal{N} \mu](x). 
\end{align*}

In the case that the measure is supported on finite number of atoms, it can be written as
\begin{align*}
	[\mathcal{N}\mu](x) = \sum_{ n = 1}^N c_n \sigma(a_n \cdot x + b_n),
\end{align*}
which is the representation of the neural network, and the gradient of the model is:
\begin{align*}
	[\partial_{x_i} \mathcal{N}](x) = \sum_{n = 1}^N c_n \partial_{x_i} \sigma(a_n \cdot x + b_n).
\end{align*}

Next we need to incorporate the nonconvex penalty in the integral representation. 
For the atom parts it is defined in the discrete setting. 
For the continuous part on $\Omega$ the measure is penalized with the weight $\phi'(0)$. 
Under assumption \ref{pen1} it is just the $L^1$ norm. 
Concretely, observe $\Omega = \atom \mu \cup \cont\mu $, we define the penalty term
\begin{equation}
	\Phi_1(\mu) = \|\mu_{\cont}\|_{M(\Omega)} + \sum_{\omega \in \atom \mu} \phi(|\mu|(\omega)).
\end{equation}

We formulate now the optimization problem as following:
\begin{equation}
	\label{int_noncon}
	\tag{$P_{\Phi_1}$}
	\left\{
	\begin{aligned}
	& \min_{\mu \in M(\Omega)} J (\mu):= L(\mu) + \alpha \Phi_1(\mu),\\
	& \text{with},\\
	& L(\mu) = \frac{1}{2} \bigg( \|[\mathcal{N}\mu] - V \|^2_{L^2(D, \nu)} + \sum_{i = 1}^d \|  \partial_{x_i}[\mathcal{N}\mu] - \partial_{x_i} V \|_{L^2(D, \nu)}^2 \bigg) .
	\end{aligned}
	\right.
\end{equation}



First, we observe the optimization problem with $L^2$-constrained inner weights and outer weights. 
We prove that in this setting the nice property of $p$-homogeneous active function still holds: 
we can constrain the inner weights on a unit ball and only penalize the outerweight with proper exponent.


\begin{proposition}[Generalization of {\cite[Proposition~3]{pieper2022nonconvex}} to $m$-homogeneous activations]
	Let $p \geq 1$, $q \geq 1$, and let $r \colon \mathbb{R}^{d+1} \to \mathbb{R}_+$ be a positively $1$-homogeneous functional:
	\[
	  r(\tau\,\omega) = |\tau|\, r(\omega) \quad \text{for all } \omega \in \mathbb{R}^{d+1},\; \tau \in \mathbb{R}.
	\]
	Suppose the activation function $\sigma(x;\omega)$ is positively $m$-homogeneous in $\omega$, i.e.,
	\[
	  \sigma(x;\, \tau\,\omega) = \tau^m\, \sigma(x;\omega) \quad \text{for all } \tau > 0,\; \omega \in \mathbb{R}^{d+1},\; x \in D.
	\]
	Then the unconstrained problem
	\begin{equation*}
	  \min_{N \in \mathbb{N},\; (c_n, \omega_n) \in \mathbb{R} \times \mathbb{R}^{d+1}}
	  \; l(\mathcal{N}_{\omega,c};\, y)
	  + \alpha \sum_{n=1}^{N} \left[ \frac{1}{p}\,|c_n|^p + \frac{1}{q}\,r(\omega_n)^q \right],
	\end{equation*}
	and the sphere-constrained problem
	\begin{equation*}
	  \min_{N \in \mathbb{N},\; \{c_n\} \in \mathbb{R},\; \{\tilde\omega_n:\, r(\tilde\omega_n) \leq 1\}}
	  \; l(\mathcal{N}_{\tilde\omega,c};\, y)
	  + \frac{2\alpha}{s}\, m^{-mp/(mp+q)} \sum_{n=1}^{N} |c_n|^{s/2},
	\end{equation*}
	where
	\[
	  s \;=\; \frac{2pq}{mp + q},
	\]
	are equivalent.
\end{proposition}
  
\begin{proof}
	Since $\sigma$ is positively $m$-homogeneous, for any $\tau > 0$ we have
	\[
		c_n\,\sigma(x;\,\omega_n)
		= c_n\,\sigma(x;\,\tau_n\,\tilde\omega_n)
		= c_n\,\tau_n^m\,\sigma(x;\,\tilde\omega_n),
	\]
	so $\mathcal{N}_{\omega,c} = \mathcal{N}_{\tilde\omega,\tilde c}$ where $\tilde\omega_n = \omega_n / \tau_n$, $\tilde c_n = c_n\,\tau_n^m$, for any choice of $\tau_n > 0$. In particular, the loss term $l(\mathcal{N}_{\omega,c};\,y)$ is
	invariant under this reparametrization. 
	For each $\omega_n \neq 0$, write $\omega_n = \tau_n\,\tilde\omega_n$ with $\tau_n = r(\omega_n) > 0$ and $r(\tilde\omega_n) = 1$ (using the $1$-homogeneity of $r$). 
	Then unconstrained problem is equivalent to
	\[
	\min_{\substack{N \in \mathbb{N},\; \{c_n\} \in \mathbb{R}^N \\
	\{\tilde\omega_n:\, r(\tilde\omega_n) \leq 1\},\; \tau_n \geq 0}}
	\; l(\mathcal{N}_{\tilde\omega,\tilde c};\, y)
	+ \alpha \sum_{n=1}^{N}
	\left[ \frac{1}{p}\,\frac{|\tilde c_n|^p}{\tau_n^{mp}} + \frac{1}{q}\,\tau_n^q \right],
\]
\quad\\ 
For fixed $\tilde c_n \neq 0$, define
\[
	g(\tau) \;=\; \frac{1}{p}\,|\tilde c_n|^p\,\tau^{-mp} + \frac{1}{q}\,\tau^q.
\]
Since $g(\tau) \to +\infty$ as $\tau \to 0^+$ and $\tau \to +\infty$, the minimum exists and is unique. Differentiating:
\[
	g'(\tau) = -\,m\,|\tilde c_n|^p\,\tau^{-mp-1} + \tau^{q-1} = 0
	\quad\Longleftrightarrow\quad
	\tau^{mp+q} = m\,|\tilde c_n|^p.
\]
Hence
\begin{equation*}
	\tau_n^* = \bigl(m\,|\tilde c_n|^p\bigr)^{1/(mp+q)}.
\end{equation*}

\quad\\
Substituting $\tau_n^*$ into $g$:

\emph{First term:}
\[
	\frac{1}{p}\,|\tilde c_n|^p \cdot \bigl(m\,|\tilde c_n|^p\bigr)^{-mp/(mp+q)}
	= \frac{m^{-mp/(mp+q)}}{p}\;|\tilde c_n|^{p\,-\,mp^2/(mp+q)}
	= \frac{m^{-mp/(mp+q)}}{p}\;|\tilde c_n|^{pq/(mp+q)}.
\]
Here we used $p - mp^2/(mp+q) = p(mp+q-mp)/(mp+q) = pq/(mp+q)$.

\emph{Second term:}
\[
	\frac{1}{q}\,\bigl(m\,|\tilde c_n|^p\bigr)^{q/(mp+q)}
	= \frac{m^{q/(mp+q)}}{q}\;|\tilde c_n|^{pq/(mp+q)}.
\]
Adding the two terms and letting $k = pq/(mp+q) = s/2$:
\[
	g(\tau_n^*)
	= \left(\frac{m^{-mp/(mp+q)}}{p} + \frac{m^{q/(mp+q)}}{q}\right) |\tilde c_n|^{k}.
\]
We simplify the coefficient. Factor out $m^{q/(mp+q)}$:
\[
	\frac{m^{-mp/(mp+q)}}{p} + \frac{m^{q/(mp+q)}}{q}
	= m^{q/(mp+q)}\!\left(\frac{m^{-1}}{p} + \frac{1}{q}\right)
	= m^{q/(mp+q)}\,\frac{q + mp}{mpq}.
\]
Since $s = 2pq/(mp+q)$, we have $(mp+q)/(mpq) = 2/(ms)$, so
\[
	g(\tau_n^*) = \frac{2}{s}\,m^{q/(mp+q)}\,m^{-1}\;|\tilde c_n|^{s/2}
	= \frac{2}{s}\,m^{q/(mp+q) - 1}\;|\tilde c_n|^{s/2}
	= \frac{2}{s}\,m^{-mp/(mp+q)}\;|\tilde c_n|^{s/2}.
\]
(using $q/(mp+q) - 1 = -mp/(mp+q)$).
\end{proof}

\begin{remark}\leavevmode
	For $p = q = 2$, the exponent becomes $s/2 = 2/(m+1)$, and the constant is $(m+1)\,m^{-m/(m+1)}$, so the sphere penalty is
	\[
	\alpha\,(m+1)\,m^{-m/(m+1)} \sum_{n=1}^N |c_n|^{2/(m+1)}.
	\]
	For $m = 1$ this gives $2\sum|c_n|$; for $m = 2$ this gives $3 \cdot 2^{-2/3}\sum|c_n|^{2/3}$.
\end{remark}

For enforcing the sparsity we implement the explicit sparsity-promoting penalty term. Recently, there is the thought-provoking result in \cite{zhang2021understanding} that indicates normal explicit generalization methods including weight decay do not enhance the generalization performance for large neural network. Nevertheless, the explicit regularization still plays a central role in the sparse learning. The network training considered here is based on the following optimization problem:
\begin{equation}
\label{opt_dis}
\begin{aligned}
		\min_{N \in \mathbb{N}, \|a_n\|^2 + |b_n|^2 \le  1 , c_n\in \mathbb{R} , n = 1, \dots, N.} l_K(\mathcal{N}_{\omega, c}(x) ;V) + \alpha \sum_{n = 1}^N \phi(|c_n|),
\end{aligned}
\end{equation}
 where the penalty term $\phi$ is observed under the following Assumption \ref{pen2}, \ref{pen3} .
 
The following proposition states that, any optimal solution of (\ref{opt_dis}) can be contracted on the unit sphere without increasing its value of the cost function. Hence we may constrain the optimization problem to a unit sphere. 


\begin{proposition}[homogeneous reformulation]
Assuming (\ref{opt_dis}) admits at least 1 global optimal solution and $\sigma$ is $p$-homogeneous with $p \ge 1$, then (\ref{opt_dis}) is equivalent to 
\begin{equation}
\label{discrete nonconvex problem}
\tag{$P_\phi$}
		\min_{N \in \mathbb{N}, \|a_n\|^2 + |b_n|^2 =  1 , c_n\in \mathbb{R} , n = 1, \dots, N.} l_K(\mathcal{N}_{\omega, c}(x) ;V) + \alpha \sum_{n = 1}^N \phi(|c_n|).
\end{equation}
\end{proposition}
\begin{proof}
    By assumption the optimal solution exits, denote $(\hat{N}, (\hat{a}_n, \hat{b}_n, \hat{c}_n)_{n = 1}^{\hat{N}}$) one of the global minimizer. The optimal value of the first part in the fidelity term is 
	\begin{align*}
		 \frac{1}{2K} \sum_{k = 1}^K \bigg|\bigg( \sum_{n = 1}^{\hat{N}} \hat{c}_n \sigma(\hat{a}_n \cdot x_k + \hat{b}_n) \bigg) - V(x_k) \bigg|^2.
	\end{align*}
    For $(\hat{a}_n, \hat{b}_n)$ in the interior of the unit ball with $\|\hat{a}_n\|^2 + |\hat{b}_n|^2 \lneqq  1$, denote $\lambda_n =\frac{1}{ \sqrt{ \|\hat{a}_n\|^2 + |\hat{b}_n|^2} }\gneqq  1$. Projecting $(\hat{a}_n, \hat{b}_n)$ on the sphere by defining
        \begin{align*}
            (\tilde{a}_n,\tilde{b}_n) = \lambda_n (\hat{a}_n, \hat{b}_n).
        \end{align*}
    Rescaling the outer weights by define $\tilde{c}_n = \frac{1}{\lambda_n^p} \hat{c}_n$, by $p$-homogeneity of the activation the function value stays the same:  
	\begin{align*}
		\frac{1}{2K} \sum_{k = 1}^K \bigg|\bigg( \sum_{n = 1}^{\hat{N}} \tilde{c}_n \sigma(\tilde{a}_n \cdot x_k + \tilde{b}_n) \bigg) - V(x_k) \bigg|^2 
		&= \frac{1}{2K} \sum_{k = 1}^K \bigg|\bigg(\sum_{n = 1}^{\hat{N}} \hat{c}_n \frac{1}{\lambda_n^p}  \lambda_n^p \sigma(\hat{a}_n \cdot x_k + \hat{b}_n) \bigg) - V(x_k) \bigg|^2, \\
		&= \frac{1}{2K} \sum_{k = 1}^K \bigg|\bigg( \sum_{n = 1}^{\hat{N}} \hat{c}_n \sigma(\hat{a}_n \cdot x_k + \hat{b}_n) \bigg) - V(x_k) \bigg|^2.
	\end{align*}
    The fidelity is invariant w.r.t. the scaling. The similar result holds for the second part of the fidelity term by linearity
	\begin{align*}
		\tilde{c}_n \nabla_x \sigma (\tilde{a}_n \cdot x_k + \tilde{b}_n) = \hat{c}_n \frac{1}{\lambda_n^p} \nabla_x \bigg(\lambda_n^p \sigma(\hat{a}_n \cdot x+ \hat{b}_n) \bigg) = \hat{c}_n \nabla_x \sigma(\hat{a}_n \cdot x + \hat{b}_n).
	\end{align*} 
	
    On the other hand, the penalty term is nondecreasing by Assumption \ref{pen2}. Since $\lambda_n > 1$ and $p \ge 1$, it holds $|\tilde{c}_n| = |\hat{c}_n|/\lambda_n^p \le |\hat{c}_n|$, hence
	\begin{align*}
		\phi(|\tilde{c}_n|) \le \phi(|\hat{c}_n|),
	\end{align*}
	The pair $\bigg(\tilde{N}, (\tilde{a}_n, \tilde{b}_n, \tilde{c}_n) \bigg)_{n = 1}^{\tilde{N}}$ with
	\begin{align*}
		\tilde{N} = \hat{N} \in \mathbb{N},\quad (\tilde{a}_n,\tilde{b}_n) \in \mathbb{S}^d, \quad \tilde{c}_n \in \mathbb{R}, \quad n = 1, \dots, \tilde{N},
	\end{align*}
	 admits smaller value. The statement holds. 
\end{proof}

% For any $\omega \neq 0$ in $\mathbb{R}^{d+1}$ (not just in a unit ball), we can always write:

% $$c \cdot \sigma(x;, \omega) = c \cdot |\omega|^p \cdot \sigma!\left(x;, \frac{\omega}{|\omega|}\right) = \tilde{c} \cdot \sigma(x;, \tilde{\omega})$$

% where $\tilde{\omega} = \omega/|\omega| \in \mathbb{S}^d$ and $\tilde{c} = c|\omega|^p$. This works for any starting $\omega$, regardless of its norm — inside a ball, outside a
% ball, on a ball of radius $R$, etc. The norm gets absorbed into the outer weight.

The proof relies essentially on the homogeneity of the activation function and the monotonicity of the penalty functions.


\section{existence of minimizers}
In this section we prove the existence of the global minimizer. We leverage the direct method of calculus variation. Since we are working on the space of measure $M(\Omega)$ we consider the weak -* convergence. As usual, the common difficulty is to prove the weak-* l.s.c.. Furthermore, the coercivity is compromised by introducing the nonconvex penalty term $\Phi$. The next lemma address this problem. It  proves that the coercivity of $\Phi$ is preserved by the Assumption \ref{pen1} of $\phi$.  

\begin{lemma}[On the coercivity of $\Phi$] \label{envolop}
Under the Assumption \ref{pen1}, for any $\mu \in M(\Omega)$, it holds
\begin{align*}
	\phi(\|\mu\|_{M(\Omega)}) \le \Phi(\mu) \le \|\mu\|_{M(\Omega)}. 
\end{align*}
In particular, the functional $\Phi$ is coercive. 
\end{lemma}
\begin{proof}
The proof relies on the subadditivity and $\phi(|z|) \le |z|$. 
\begin{align*}
	\phi(\|\mu\|_{M(\Omega)}) &= \phi \bigg(|\mu|(\Omega \backslash \atom \mu) + \sum_{\omega \in \atom \mu } |\mu|(\{\omega\}) \bigg),\\
	(\text{subadditivity) } 
	& \le \phi(|\mu|(\Omega\backslash \atom \mu) + \sum_{\omega \in \atom \mu } \phi(|\mu|(\{\omega\})), \\
	\bigg(\phi(|z|) \le |z|\bigg)\quad
	& \le |\mu|(\Omega\backslash \atom \mu) + \sum_{\omega \in  \atom \mu} \phi(|\mu|(\{\omega\})),\\
	& = \Phi(\mu).
\end{align*}
In particular, by the coercivity of $\phi$ it holds that $\Phi$ is also coercive. \\
For the second inequality, we leverage that $\phi$ is  sublinear and $\phi(0)= 0$. We get $\Phi(\mu) \le \|\mu\|_{M(\Omega)}$.  
\end{proof}

The next lemma proves the weak-* l.s.c. of $\Phi$. 
The proof follows the idea in the paper \cite{bouchitte1990new}. 
Since the original proof is valid for $\Omega$ locally compact, the result applies to the current setting. 
Moreover, since our functional doesn't have the singular part, the proof and assumptions can be simplified. 
\begin{lemma}
Given $\mu = \mu_{\cont} + \mu_{\atom}$, under Assumption \ref{pen1}, \ref{pen2}, it holds 
\begin{align*}
	\Phi(\mu) = \int_\Omega d\mu_{\cont} + \sum_{ \omega \in \atom \mu } \phi(|\mu_{\atom}|(\{\omega\}))
\end{align*}
is weak-* sequentially l.s.c..
\end{lemma}
\begin{proof} 
	By proposition 2.5 in \cite{bouchitte1990new} it is equivalent to show the weak-* l.s.c. of $\Phi_\epsilon(\mu): = \Phi(\mu) + c \|\mu\|_{M(\Omega)}$ for each $c > 0$.
	We can therefore assume there exists $c > 0$, such that 
	\begin{align*}
		\Phi(\mu) \ge c |\mu|(\Omega) \quad \forall \mu \in M(\Omega).
	\end{align*}
	Let $\mu_h \overset{*}{\rightharpoonup} \mu$, we decompose $\Omega$ by heavy atomes and small atoms combined the continuous part:
	\begin{align*}
		\Omega = A_h + (\Omega - A_h), \text{ where } A_h:= \{ \omega \in \Omega: \mu_h(\omega) \ge \epsilon \}.
	\end{align*}
	By passing to a subsequence we have 
	\begin{equation*}
		\mu_{h, \cont} \overset{*}{\rightharpoonup} \mu_{\cont}, \quad 
		\mu_{h, \atom} \overset{*}{\rightharpoonup} \mu_{\atom}
	\end{equation*}
	We prove the weak-* l.s.c. of each part. \\
	\textbf{On $ A_h $:} \par   
	Since $\mu_h$ has bounded variation and $\mu_h(\omega) \ge \epsilon$ for $\omega \in A_h$, we know the cardinality of $A_h$ is bounded. 
	By lemma 3.6 in \cite{bouchitte1990new} we know $M(A_h)$ is sequentially weak-* closed.
	Moreover, $\Phi(\mu_h \mathbbm{1}_{A_h})$ is weak-* l.s.c. if $\phi$ is l.s.c. and subadditive, therefore we have
	\begin{align*}
		\liminf_{h} \Phi(\mu_h \mathbbm{1}_{A_h})  
		= \liminf  \sum_{ \omega \in \atom \mu_h \cap A_h } \phi(|\mu_{h, \atom}|(\{\omega\}))
		\ge \Phi ( \mu \mathbbm{1}_{A_h}).
	\end{align*}
	
	\quad\\
	\textbf{On $\Omega - A_h$:} \par   
	By proposition \ref{secondorderestimate}, we have for $\epsilon$ small enough, for all $z \in [0, \epsilon]$: 
	\begin{equation*}
		\begin{aligned}
			&\phi(z) \ge z - \frac{\gamma_1 z^2}{2}, \\
			&\phi(z) \le z.
		\end{aligned}
	\end{equation*}
	We have
	\begin{equation*}
	\begin{aligned}
		&\liminf_h \Phi(\mu_h \mathbbm{1}_{\Omega - A_h}) \\
		= & \liminf_h \bigg( \int_{\Omega - A_h} d\mu_{\cont} + \sum_{\omega \in (\Omega - A_h) \cap \atom \mu_h} \phi(|\mu_{h, \atom}|(\{\omega\})) \bigg)\\
		\ge & \liminf_h \bigg( \int_{\Omega - A_h} d\mu_{\cont} + \sum_{\omega \in (\Omega - A_h) \cap \atom \mu_h} |\mu_{h, \atom}|(\{\omega\})  - \frac{\gamma_1}{2} |\mu_{h, \atom}|(\{\omega\})^2  \bigg).
	\end{aligned}
	\end{equation*}
	Choose $\delta \ge \frac{\gamma_1}{2} \epsilon$, we have
	\begin{equation*}
		\sum_{\omega \in (\Omega - A_h) \cap \atom \mu_h} |\mu_{h, \atom}|(\{\omega\})  - \frac{\gamma_1}{2} |\mu_{h, \atom}|(\{\omega\})^2 
		\ge \sum_{\omega \in (\Omega - A_h) \cap \atom \mu_h} ( 1 - \delta) |\mu_{h, \atom}|(\{\omega\}).
	\end{equation*}
	Hence
	\begin{equation*}
		\begin{aligned}
			&\liminf_h \Phi(\mu_h \mathbbm{1}_{\Omega - A_h}) \\
			\ge & \liminf_h (1 - \delta) \int_{\Omega - A_h} d\mu_{\cont} + \liminf_h( 1 - \delta) \sum_{\omega \in (\Omega - A_h) \cap \atom \mu_h} |\mu_{h, \atom}|(\{\omega\})\\
			= & \liminf_h (1 - \delta) \|\mu_h \mathbbm{1}_{\Omega - A_h}\|_{M(\Omega)}.
		\end{aligned}
	\end{equation*}
	By the weak-* l.s.c. of the norm $\|\cdot\|_{M(\Omega)}$ we have 
	\begin{equation*}
		\liminf_h \Phi(\mu_h \mathbbm{1}_{\Omega - A_h}) \ge (1 - \delta) \|\mu \mathbbm{1}_{\Omega - A_h}\|_{M(\Omega)}.
	\end{equation*}
	Furthermore, applying $\phi(z) \le z$ we get
	\begin{align*}
		\sum_{\omega \in (\Omega - A_h) \cap \atom \mu_h}  |\mu_{\atom}|(\{\omega\})
		\ge \sum_{\omega \in (\Omega - A_h) \cap \atom \mu_h} \phi(|\mu_{h, \atom}|(\{\omega\})).
	\end{align*}
	Therefore
	\begin{equation*}
		\begin{aligned}
			& \liminf_h \Phi(\mu_h \mathbbm{1}_{\Omega - A_h})\\
			\ge & (1 - \delta) \|\mu \mathbbm{1}_{\Omega - A_h}\|_{M(\Omega)}\\
			\ge & (1 - \delta)\int_{\Omega - A_h} d\mu_{\cont} + \sum_{\omega \in (\Omega - A_h) \cap \atom \mu_h} \phi(|\mu_{h, \atom}|(\{\omega\}))\\
			= & (1 - \delta) \Phi(\mu \mathbbm{1}_{\Omega - A_h}).
		\end{aligned}
	\end{equation*}
Taking $\epsilon \rightarrow 0$ we get $\delta \rightarrow 0$ and the result is proven. 
\end{proof}

\begin{theorem}[Existence of the global minimizer]
Assuming Assumption \ref{pen1}, \ref{pen2} is fulfilled, problem (\ref{int_noncon}) admits at least one global minimizer in $M(\Omega)$. 
\end{theorem}
\begin{proof}
We apply the direct method of the calculus of variation:\\ 
1. Boundedness (ability to take a minimizing sequence in $M(\Omega)$): $J \ge 0$, therefore $J$ is bounded from below.\\
2. Minimizing sequence admits a weak-* convergent subsequence: To apply Banach-Alaoglu we need to prove that every minimizing sequence of $J$ is bounded in $M(\Omega)$.  We prove it by the usual coercivity argument. By the non-convexity of the penalty term, the coercivity of the cost function is compromised. However, we get form lemma \ref{envolop}
\begin{align*}
\alpha\|\mu\|_{M(\Omega)} \le \alpha \Phi(\mu) \le J(\mu),
\end{align*}
hence $J$ is also coercive.  \\
3. Lower-semi continuity: by the lemma $\Phi(\mu)$ is weak-* l.s.c., therefore we have for $\mu_n \overset{*}{\rightharpoonup} \hat{\mu}$, 
\begin{align*}
\liminf_{n \rightarrow \infty } J(\mu_n) &= \liminf_{n \rightarrow \infty} \bigg( L(\mu_n)  + \alpha \Phi(\mu_n) \bigg)\\
&\ge L(\hat{\mu}) + \alpha \liminf_{n \rightarrow \infty} \Phi(\mu_n)\\
& \ge L(\hat{\mu}) + \alpha\Phi(\hat{\mu})\\
& = J(\hat{\mu}),
\end{align*}
where the inequality in the second line follows from the weak-* l.s.c. of $L$ (since $L$ is convex and continuous composed with the weak-*-to-weak continuous operator $\mathcal{N}$) and the inequality in the third line follows from weak-* l.s.c. of the functional $\Phi$.
\end{proof}
 
\section{optimality condition and finite support of the solution}
Since the optimization problem is nonconvex, the global minimizer is difficult to find in general. 
Instead, we define the local minimum: 
\begin{definition}[local minimizer]
$\bar{\mu} \in M(\Omega)$ is a local minimizer if there exists $\epsilon > 0$ such that 
\begin{align*}
	J(\bar{\mu}) \le J(\mu) \quad \text{for all } \mu \text{ with }  \|\mu - \bar{\mu}\|_{M(\Omega)} \le \epsilon. 
\end{align*}
\end{definition}

\begin{remark}[The gradient of $L$]
To further investigate the optimality condition of (\ref{int_noncon}) and derive the according algorithm, we calculate the first variation of $L(\mu)$,  We define the gradient information $\bar{p}(\omega) \in C(\Omega)$ as
 \begin{equation}
 \begin{aligned}
 	\bar{p}(\omega) &:= \nabla L(\bar{\mu}) \\
 	& =  \mathcal{N}^*(\mathcal{N}[\bar{\mu}] - V) + \sum_{i = 1}^d (\partial_{x_i} \mathcal{N})^*( \partial_{x_i} \mathcal{N}[\bar{\mu}] - \partial_{x_i} V ) \\
 	& = \int_D \sigma(x, \omega) (\mathcal{N}[\bar{\mu}](x) - V) d\nu(x) + \sum_{i = 1}^d \int_D \partial_{x_i} \sigma(x, \omega) ( \partial_{x_i} \mathcal{N}[\bar{\mu}](x) - \partial_{x_i}V ) d\nu(x) .
 	 \end{aligned}
 \end{equation}
 For the finite sample and $\sigma(x) = \relu(x)^p$, the gradient admits the form:
 \begin{equation}
 	p(\omega) = \frac{1}{K} \sum_{k = 1}^K  \relu(\omega \cdot x_k)^p (\mathcal{N}_{\bar{\mu}}(x_k) - y_k) + 
 	\frac{1}{K} \sum_{k = 1}^K \sum_{i = 1}^d p \cdot \relu( \omega \cdot x_k )^{ p -1}  \omega_i (\partial_{x_i} \mathcal{N}(x_k) - \partial_{x_i}V(x_k))
 \end{equation}
\end{remark} 
 
 
\begin{theorem}[the optimality condition]
\label{opt_cdt}
	Under the Assumption \ref{pen1}, \ref{pen2}, if $\bar{\mu}$ is a local minimum of $J$, then it holds
	\begin{enumerate}
		\item $|\bar{p}| \le \alpha \qquad \text{for all } \omega$.
		\item $\bar{p}(\omega) = -\alpha \phi'(|\bar{\mu}|(\omega)) \sign(\bar{\mu}\{ w \}), \qquad \text{for } \bar{\mu}\text{-almost all } \omega \in \Omega.$
	\end{enumerate}
	the signum function $\sign: \Omega \rightarrow \{-1, 1\}$ is defined $\mu$-a.s. by Hahn decomposition. 
\end{theorem}
\begin{proof} 
We leverage that dual evaluation of $\delta_\omega$ is the evaluation of $\bar{p}$ at $\omega$: $\langle \bar{p}, \delta_\omega\rangle_{C(\Omega), M(\Omega)}= \bar{p}(\omega)$. Since $\bar{\mu}$ is the local minimizer, it holds for all $u \in M(\Omega)$:
\begin{align*}
	J(\bar{\mu} + \tau u) - J(\bar{\mu}) = L(\bar{\mu} + \tau u) - L(\bar{\mu}) + \alpha \Phi(\bar{\mu} + \tau u) - \alpha \Phi(\bar{\mu}) \ge 0.
\end{align*}
Taking $\tau \rightarrow 0$, it holds
\begin{align*}
	-\langle \bar{p}, u\rangle \le \alpha \lim_{ \tau \rightarrow 0} \frac{1}{\tau} (\Phi(\bar{\mu} + \tau u) - \Phi(\bar{\mu}) ). 
\end{align*}
Now we separate 2 cases:
\quad\\
\textbf{1)} For the atom part,  $\omega \in \atom \bar{\mu} $ with $\bar{\mu}(\{\omega\}) = c$, taking $u = \delta_\omega$.
it holds
\begin{align*}
	- p(\omega) =  -\langle \bar{p}, u \rangle \le \alpha \lim_{\tau \rightarrow 0} \frac{1}{\tau} (\phi (|c + \tau|) - \phi(|c|) = \alpha \phi'(|c|) \sign(c).  
\end{align*}
Taking $u = \pm \delta_\omega$, we get $p(\omega) = -\alpha \phi'(|c|) \sign(c)$. 
Since $\phi'(0) = 1$ and $\phi'$ is decreasing we have $\phi'(|c|) \in [0, 1]$, hence $|\bar{p}| \le \alpha$. 

\quad\\
\textbf{2)} For the continuous part, first observe $\omega \notin \atom \bar{\mu}$ and perturb with $ u = \delta_\omega$, we have
\begin{align*}
	-\bar{p}(\omega) =  - \langle p, u \rangle \le \alpha \lim_{\tau \rightarrow 0 } \frac{1}{\tau} ( \phi(\tau) - \phi(0) ) = \alpha \phi'(0) = \alpha.
\end{align*}
Taking $u = - \delta_\omega$ we get $|p(\omega)| \le \alpha$ for all $\omega \notin \atom \bar{\mu}$.
Next we take $u = \bar{\mu}_{\text{cont}}$, we have
\begin{align*}
	-\langle \bar{p}, \bar{\mu}_{\text{cont}} \rangle \le \alpha \lim_{\tau \rightarrow 0} \frac{1}{\tau} \tau  \|\bar{\mu}_{\text{cont}}\|_{M(\Omega)} =  \alpha \|\bar{\mu}_{\text{cont}}\|_{M(\Omega)}.
\end{align*}
Taking $u = - \bar{\mu}_{\text{cont}}$ we get $-\langle \bar{p}, \bar{\mu}_{\text{cont}} \rangle  = \alpha \|\bar{\mu}_{\text{cont}}\|_{M(\Omega)}$, i.e., 
\begin{align*}
	\int_\Omega \bigg( - \bar{p}(\omega) * \sign \bar{\mu}_{\cont} (\omega)  - \alpha \bigg) d|\bar{\mu}|(\omega) = 0. 
\end{align*}
By $|p(\omega)| \le \alpha$ we get $\bar{p} =  - \alpha *  \sign \bar{\mu}_{\cont}$ for almost all $\omega \in \cont \bar{\mu}$. 
\end{proof}

The optimality condition assures the existence of an optimizer. 
In absence of Assumption \ref{pen1}, it is difficult to prove that a global optimizer exists. 
For example for the penalty $ \phi(c) = |c|^q$, we have $\phi'(0) = \infty$, the weak-* l.s.c. of the cost functional breaks.
However, assuming the local optimizer exists, we can still derive the necessary optimality condition. 

\begin{corollary}[Optimality conditions for $\phi'(0^+) = +\infty$]                     
\label{cor:singular_penalty}                                                                                                                                                                                                                
Let $\Omega$ be compact and suppose $\phi$ satisfies Assumption~\ref{pen2} but with $\phi'(0^+) = +\infty$ in place of Assumption~\ref{pen1}                                                                                                
(e.g., $\phi(z) = \frac{1}{q}z^q$ for $0 < q < 1$).                                                                                                                                                                                         
If $\bar{\mu}$ is a local minimizer of \eqref{int_noncon} with atomic part                                                                                                                                                                  
$\bar{\mu}_{\atom} = \sum_{n=1}^N \bar{c}_n\,\delta_{\bar{\omega}_n}$, $\bar{c}_n \neq 0$, then:
\begin{enumerate}
	\item Stationarity at each atom:
	\[
		\bar{p}(\bar{\omega}_n) = -\alpha\,\phi'(|\bar{c}_n|)\,\sign(\bar{c}_n),
		\qquad n = 1, \dots, N.
	\]

	\item Lower bound on atom weights:
	Since $\bar{p} \in C(\Omega)$ and $\Omega$ is compact,
	$\|\bar{p}\|_{C(\Omega)} < \infty$.
	Combined with (1), this gives
	$\phi'(|\bar{c}_n|) = |\bar{p}(\bar{\omega}_n)|/\alpha \le \|\bar{p}\|_{C(\Omega)}/\alpha$
	for every $n$. Since $\phi'$ is decreasing with $\phi'(0^+) = +\infty$,
	inverting yields a uniform lower bound:
	\[
		|\bar{c}_n| \;\ge\; (\phi')^{-1}\!\Big(\frac{\|\bar{p}\|_{C(\Omega)}}{\alpha}\Big)
		\;>\; 0,
		\qquad n = 1, \dots, N.
	\]
	In particular, $\bar{\mu}$ has at most finitely many atoms:
	$N \le \|\bar{\mu}\|_{M(\Omega)} \big/ \min_n |\bar{c}_n| < \infty$.
\end{enumerate}
\end{corollary}
\begin{proof}
(1) follows by the same directional derivative argument as in Theorem~\ref{opt_cdt}.
For $\omega_n \in \atom\bar{\mu}$ with $\bar{\mu}(\{\omega_n\}) = \bar{c}_n \neq 0$,
perturbing with $u = \pm\delta_{\omega_n}$ gives
\[
	\mp\,\bar{p}(\omega_n)
	\;\le\; \alpha\,\lim_{\tau \to 0}\frac{1}{\tau}
	\bigl[\phi(|\bar{c}_n \pm \tau|) - \phi(|\bar{c}_n|)\bigr]
	\;=\; \pm\,\alpha\,\phi'(|\bar{c}_n|)\,\sign(\bar{c}_n),
\]
which yields equality. Note that $\phi'(|\bar{c}_n|) < \infty$ for $\bar{c}_n \neq 0$,
so this step does not require $\phi'(0) < \infty$.

\medskip\noindent
(2): The dual variable $\bar{p} = \nabla L(\bar{\mu}) \in C(\Omega)$
is continuous on the compact set $\Omega$, hence bounded:
$\|\bar{p}\|_{C(\Omega)} < \infty$.
By (1), $\phi'(|\bar{c}_n|) = |\bar{p}(\bar{\omega}_n)|/\alpha
\le \|\bar{p}\|_{C(\Omega)}/\alpha$.
Since $\phi'$ is strictly decreasing (by concavity of $\phi$) and
$\phi'(0^+) = +\infty$, the function $\phi'$ restricted to $(0,\infty)$
has a well-defined decreasing inverse, and
\[
	|\bar{c}_n| \;\ge\; (\phi')^{-1}\!\Big(\frac{\|\bar{p}\|_{C(\Omega)}}{\alpha}\Big)
	\;=:\; c_{\min} \;>\; 0.
\]
The finiteness of the support then follows from the bounded total variation:
$N \cdot c_{\min} \le \sum_{n=1}^N |\bar{c}_n| \le \|\bar{\mu}\|_{M(\Omega)} < \infty$.
\end{proof}

\begin{remark}
Compared with Theorem~\ref{opt_cdt}:
\begin{enumerate}
	\item The global dual bound $|\bar{p}(\omega)| \le \alpha$ for all $\omega \in \Omega$
	is \emph{lost}: for $\omega \notin \atom\bar{\mu}$, perturbing with
	$u = \delta_\omega$ yields
	$-\bar{p}(\omega) \le \alpha\,\phi'(0^+) = +\infty$, which is vacuous.
	No constraint on $\bar{p}$ is obtained away from the support.
	\item In exchange, the singularity $\phi'(0^+) = +\infty$ provides an
	\emph{automatic minimum atom size} via (2), and finite support
	follows without the merging argument of Theorem~5.2 or
	the local strong concavity in Assumption~\ref{pen3}.
	\item The zero measure $\mu = 0$ is always a local minimizer,
	since adding any atom with small weight $|c| \ll 1$ incurs a
	marginal penalty $\phi(|c|)/|c| \to +\infty$ that dominates the
	potential decrease in the loss.
	This means the node insertion criterion $|\bar{p}(\omega)| > \alpha$
	of the greedy algorithm no longer applies.
\end{enumerate}
\end{remark}

\begin{corollary}[Existence of minimizer for $\eqref{discrete nonconvex problem}$ with $\phi'(0^+) = +\infty$]
\label{cor:discrete_existence}
Suppose $\phi$ satisfies Assumption~\ref{pen2} with $\phi'(0^+) = +\infty$, that $\sigma$ is continuous and $D$ is bounded. Then the discrete problem $\eqref{discrete nonconvex problem}$ admits a global minimizer with at most $N_{\max}$ neurons, where
\[
	N_{\max} = \frac{\|V\|_H^2}{2\,\alpha\,\phi(c_{\min})}, \qquad
	c_{\min} = (\phi')^{-1}\!\bigg(\frac{\|V\|_H\;\sup_{\omega \in \mathbb{S}^d}\|\sigma(\cdot\,;\omega)\|_H}{\alpha}\bigg),
\]
where $\|\cdot\|_H$ denotes the norm induced by the loss (combining $L^2$ and gradient terms).
\end{corollary}
\begin{proof}\quad\\
1. (Fixed $N$, existence by Weierstrass).
For fixed $N$, the objective is a real-valued function on $(\mathbb{S}^d)^N \times \mathbb{R}^N$:
\[
	J(\omega, c) = \frac{1}{2} \int_D \bigg| V(x) - \sum_{n=1}^N c_n \sigma(a_n \cdot x + b_n) \bigg|^2 + \sum_{i=1}^d \bigg| \partial_{x_i} V(x) - \sum_{n=1}^N c_n \partial_{x_i} \sigma(a_n \cdot x + b_n) \bigg|^2 d\nu + \alpha \sum_{n=1}^N \phi(|c_n|).
\]
We verify $J$ is continuous in $(\omega, c)$.
Since $\sigma$ is continuous and $\mathbb{S}^d \times \overline{D}$ is compact ($D$ bounded), both $\sigma(a \cdot x + b)$ and $\nabla_x \sigma(a \cdot x + b)$ are uniformly bounded on $\mathbb{S}^d \times \overline{D}$.
For any convergent sequence $(\omega^k, c^k) \to (\omega, c)$, the integrand converges pointwise in $x$ and is dominated by $C(|V(x)|^2 + |\nabla V(x)|^2 + 1) \in L^1(D, \nu)$.
By dominated convergence, the integral is continuous in $(\omega, c)$. Combined with the continuity of $\phi$, $J$ is continuous.

Since $\phi(z) \to \infty$ as $z \to \infty$ (Assumption~\ref{pen2}) and $J \ge 0$, any minimizing sequence satisfies $\sum \phi(|c_n|) \le J_0/\alpha$, so each $|c_n|$ is bounded. Therefore the minimizing sequence lies in $(\mathbb{S}^d)^N \times [-R, R]^N$ for some $R > 0$, which is compact. By Weierstrass, $J$ attains its minimum for each $N$.

\medskip\noindent
2. (Uniform bound on $N$).
Let $(\omega^*, c^*)$ be a minimizer for $N$ neurons with all $c_n^* \neq 0$.
Since taking $c = 0$ is feasible and the penalty is non-negative, we have
\begin{equation*}
	L(\omega^*, c^*) \le J(\omega^*, c^*) \le J(\omega, 0) = L(0) = \frac{1}{2}\|V\|_H^2,
\end{equation*}
Analogously it holds
\begin{equation*}
	\sum_{n = 1}^N \phi(|c_n^*|) \le \frac{1}{2\alpha}\|V\|_H^2
\end{equation*}
The stationarity condition in $c_n$ gives $|p(\omega_n^*)| = \alpha\,\phi'(|c_n^*|)$, where $p(\omega) = \langle \mathcal{N}_{\omega^*, c^*} - V,\, \sigma(\cdot\,;\omega)\rangle_H$.
By Cauchy--Schwarz:
\[
	\alpha\,\phi'(|c_n^*|) = |p(\omega_n^*)| \le \|\mathcal{N}_{\omega^*, c^*} - V\|_H \;\|\sigma(\cdot\,;\omega_n^*)\|_H \le \|V\|_H\;\sup_{\omega \in \mathbb{S}^d}\|\sigma(\cdot\,;\omega)\|_H.
\]
Since $\phi'$ is decreasing with $\phi'(0^+) = +\infty$ we can define $c_{\min}$ with
\begin{equation*}
	c_{\min} := (\phi')^{-1}\!\bigg(\frac{\|V\|_H\;\sup_{\omega \in \mathbb{S}^d}\|\sigma(\cdot\,;\omega)\|_H}{\alpha}\bigg),
\end{equation*} 
with $|c_n^*| \ge c_{\min} > 0$.
Hence $N\,\phi(c_{\min}) \le \sum \phi(|c_n^*|) \le \frac{1}{2\alpha}\|V\|_H^2$, giving $N \le N_{\max}$.

For $N > N_{\max}$, any $N$-neuron minimizer must have some $c_n = 0$, reducing to an $N'$-neuron solution with $N' \le N_{\max}$.
Therefore $\inf_N J_N^* = J_{N_{\max}}^*$, which is attained by Step~1.
\end{proof}

If the sample size is finite, then the optimal solution admits finite support by the Carath\'eodory's theorem. The following theorem proves that the optimal solution admits the finite support even though the dataset is infinitely large. 

\begin{theorem}
	If the Assumption \ref{pen1}, \ref{pen2}, \ref{pen3} are satisfied, then the optimal solution of (\ref{int_noncon}) admits at most finite nodes.
\end{theorem}
\begin{proof}\quad\\
1. $\cont\bar{\mu} = \emptyset$: Argue by contradiction. Assuming there exists $\hat{\omega}$ such that $\bar{\mu}(\{\hat{\omega}\}) = 0$ and $\delta > 0$ such that the open ball $B_\delta(\hat{\omega})$ on $\Omega$ satisfies $|\bar{\mu}|(B_\delta(\hat{\omega})) > 0$. 
Define $D_\delta := B_\delta(\hat{\omega})$ and $C_\delta := |\bar{\mu}|(D_\delta)$ and construct  the measure 
\begin{align*}
	\mu_\delta = \bar{\mu} - \bar{\mu}|_{D_\delta} + C_\delta \delta_{\hat{\omega}}. 
\end{align*}
We prove that the constructed measure decrease the cost function. \par
The cost function is quadratic, therefore the second order Taylor expansion is exact, we get:
\begin{align*}
	L(\mu_\delta)  &= L(\bar{\mu}) + \langle \bar{p}, \mu_\delta - \bar{\mu} \rangle_{C(\Omega), M(\Omega)} + \frac{1}{2} \|\mathcal{N}[u_\delta - \bar{\mu}]\|_{L^2(D, \nu)}^2 + \frac{1}{2}\sum_{i = 1}^d \|\partial_{x_i} \mathcal{N}[\mu_\delta - \bar{\mu}]\|_{L^2(D, \nu)}^2
\end{align*}
By Assumption \ref{assumption on actfun}, it holds
\begin{align*}
	 |\mathcal{N}[\mu_\delta - \bar{\mu}](x)| &\le \bigg|\int_{D_\delta} \sigma(x, \omega) d\bar{\mu} - C_\delta \sigma(x, \hat{\omega}) \bigg|,\\
	& = \bigg|\int_{D_\delta} \sigma(x, \omega)  - \sigma(x, \hat{\omega}) d \bar{\mu} \bigg|,\\
	& \le \delta \Lambda (1 + \|x\|)|\bar{\mu}|(D_\delta), \\
	& =  \delta \Lambda (1 + \|x\|)C_\delta. 
\end{align*}
Analogously it holds
\begin{align*}
	|\partial_{x_i} \mathcal{N}[\mu_\delta - \bar{\mu}](x)| \le \delta \Lambda (1 + \|x\|)C_\delta.
\end{align*}
Define 
\begin{align*}
	\Lambda_1^2(D, \Lambda) =  \int_D(1 + \|x\|)^2 d\nu  \Lambda^2.
\end{align*}
Hence
\begin{align*}
	L(\mu_\delta) \le L(\bar{\mu}) + \langle \bar{p}, \mu_\delta - \bar{\mu} \rangle + \frac{d + 1}{2}\delta^2 \Lambda_1^2 C_\delta^2. 
\end{align*}
By the Theorem \ref{opt_cdt} $\bar{p}$ is constant for $\omega \in \cont \bar{\mu}$, therefore the first order term vanishes, it holds
\begin{align*}
	L(\mu_\delta) \le L(\bar{\mu})  + \frac{d + 1}{2}\delta^2 \Lambda_1^2 C_\delta^2. 
\end{align*}
Concerning the cost term, it holds
\begin{align*}
	\Phi(\mu_\delta) &= \Phi(\bar{\mu}) - \bar{\mu}(D_\delta) + \phi(C_\delta),\\
	&=\Phi(\bar{\mu}) + \phi(C_\delta) - \bar{\mu}(D_\delta) , \\
	&= \Phi(\bar{\mu}) + \int_0^{C_\delta} \phi'(\xi) - \phi'(0) d\xi,\\
	\text{Assumption \ref{pen3}(2) }
	& \le \Phi(\bar{\mu}) - \frac{ \gamma_2 }{2} C_\delta^2. 
\end{align*}
Combining the 2 parts we get
\begin{align*}
	J(\mu_\delta) \le J(\bar{\mu}) - \bigg(\frac{ \alpha \gamma_2 - (d + 1) \delta^2 \Lambda_1^2}{2} \bigg) C_\delta^2.
\end{align*}
Choose $\delta^2 \le \frac{2 \alpha \gamma_2}{(d + 1)\Lambda_1^2(D, \Lambda)}$ we get the contradiction.

\quad\\
2. $|\atom \bar{\mu}|$ is finite: 
Without loss of generality we can assume the atomic support of $\mu$ is contained in a compact subset of $\Omega$. 
Assume by contradiction, there exists a sequence $\omega_n \rightarrow \infty $ escapes to the infinity.
Since $\bar{\mu} \in M(\Omega)$ admits bounded variation, by passing to a subsequence, it holds $\bar{\mu}(\omega_n) \rightarrow 0$. 
Assuming w.l.o.g. that $\bar{\mu}(\omega_n) = c_n > 0$, we have
\begin{equation*}
	\begin{aligned}
		\lim_{\omega_n \rightarrow + \infty } \bar{p}(\omega_n)
	= -\alpha \phi'(|\bar{\mu}|(\omega_n)) \sign (\bar{\mu}(\omega_n))
	= - \alpha,
	\end{aligned}
\end{equation*}
which contradicts to that $\bar{p} \in C_0(\Omega)$ and $\lim_{\omega_n \rightarrow + \infty } \bar{p}(\omega_n) \rightarrow 0$.

Now we observe $\atom \bar{\mu}$.
Since it is contained in a compact subset of $\Omega$, there exists a sequence $\omega_n$ (possibly passing through a subsequence) such that 
\begin{align*}
	\omega_n \rightarrow \hat{\omega} \text{ in $\Omega \quad$ for some } \hat{\omega} \in \Omega.
\end{align*}
Assuming w.l.o.g. that $\bar{\mu}(\omega_n) = c_n > 0$, by the boundedness of total variation it holds $ \lim_{n \rightarrow \infty} c_n = 0$. By the Theorem \ref{opt_cdt} , it holds
\begin{align*}
	\bar{p}(\hat{\omega}) = \lim_{n \rightarrow \infty} \bar{p}(\omega_n) = -\alpha \lim_{n\rightarrow \infty} \phi'(c_n) \rightarrow -\alpha.
\end{align*}
Concentrate the measure on $\{\omega_{N}, \dots\}$ on the limits $\hat{\omega}$, i.e., define
\begin{align*}
	\mu_N := \bar{\mu} - \sum_{i = N}^\infty \delta_{\omega_i} \bar{\mu}(\omega_i) + C_N \delta_{\hat{\omega}}, \text{ where } C_N = \sum_{i = N}^\infty c_i > 0.
\end{align*}
As in 1., the fidelity term of the cost can by decomposed by Taylor expansion.\\
Define $\delta_N:= \max_{i = N}^\infty |\omega_i - \hat{\omega}|$ :
\begin{align*}
	L(\mu_N) \le L(\bar{\mu}) + \langle \bar{p}, \mu_N - \bar{\mu}\rangle +\frac{d +1}{2} \Lambda_1^2(D, \Lambda) \delta_N^2  C_N^2
\end{align*}
For the first variation we have by the optimality condition and the argument above:
\begin{align*}
	\langle \bar{p}, \mu_N - \bar{\mu} \rangle = \langle \bar{p},  -\sum_{i = N}^\infty \delta_{\omega_i} \bar{\mu}(\omega_i) + C_N \delta_{\hat{\omega}} \rangle = \alpha \sum_{i = N}^\infty c_i \phi'(c_i) - \alpha C_N.
\end{align*}

For the penalty term, it holds
\begin{align*}
	\Phi(\mu_N) = \Phi(\mu) - \sum_{i = N}^\infty \phi(c_i) + \phi(C_N). 
\end{align*}
Combining the estimation, it holds
\begin{align*}
	J(\mu_N) \le J(\bar{\mu}) - \alpha \sum_{i = N}^\infty \bigg(\phi(c_i) - c_i \phi'(c_i) \bigg) - \alpha\bigg(C_N  - \phi(C_N) \bigg) + \frac{d + 1}{2} \Lambda_1^2(D, \Lambda) \delta_N^2 C_N^2.
\end{align*}
By concavity of $\phi$ it holds
\begin{align*}
	\phi(c_i) - c_i \phi'(c_i) \ge \phi(0) = 0.
\end{align*}
By Assumption \ref{pen3} (2) it holds
\begin{align*}
	C_N - \phi(C_N) &= C_N \phi'(0) - \bigg(\phi(0) + \int_0^{C_N} \phi'(\xi) d\xi \bigg), \\
	(\text{Assumption \ref{pen1}} )
	& = -\int_0^{C_N}\phi'(\xi) - \phi'(0) d\xi, \\
	& \ge \frac{\gamma_2}{2}C_N^2.
\end{align*}
Thus
\begin{align*}
	J(\mu_N) \le J(\bar{\mu}) - \frac{1}{2}\bigg(\alpha \gamma_2 -(d + 1) \Lambda_1^2(D, \Lambda)\delta_N^2 \bigg)C_N^2.
\end{align*}
Taking 
\begin{align*}
	\delta^2 < \frac{\alpha \gamma_2}{\Lambda_1^2(\Lambda, D)(d + 1)},
\end{align*}
we get the contradiction.
\end{proof}

Now we have a necessary optimality condition. But it is not sufficient. More importantly, we cannot formulate an algorithm to determine the optimal value of $N$. To overcome this difficulty, we apply a heuristic approach: For fixed $N$, the optimization problem
\begin{equation}\label{opt_fixN}
	J_{\bar{\omega}}(c):= l(\mathcal{N}_{\bar{\omega}, c}; y) + \sum_{n = 1}^N \phi(|c_n|),
\end{equation}
is well-equipped to solve by traditional gradient-based optimization method. Now we need to find a way to heuristically determine the optimal $\bar{\omega}$ and the optimal number $N$. To this end, we introduce the following theorem as a sufficient condition. 


\begin{theorem}
\label{sc_opt}
 Assuming $\phi$ satisfies the Assumption \ref{pen1}, \ref{pen2}, \ref{pen3}, let $\bar{\mu}$ be a finitely supported measure $\bar{\mu}:= \sum_{n = 1}^N \bar{c}_n \delta_{\bar{\omega}_n}$ such that 
\begin{enumerate}
	\item $\bar{c}_n \in (\mathbb{R} \backslash 0)^N$ is a local minimum of $J_{\bar{\omega}}$, 
	\item For all $\omega \in \Omega \backslash \{\bar{\omega}_n \}_{n = 1, \dots, N}$, it holds $|\bar{p}(\omega)|  \lneqq  \alpha$, where $\bar{p} = \mathcal{N}^*(\mathcal{N}_{\bar{\omega}, \bar{c}} - y) + \sum_{i=1}^d (\partial_{x_i}\mathcal{N})^*(\partial_{x_i}\mathcal{N}_{\bar{\omega}, \bar{c}} - \partial_{x_i} V)$. 
\end{enumerate}
Then $\bar{\mu}$ is a local minimum of $J$.
\end{theorem}
\begin{proof}
For arbitrary $\epsilon > 0$ we observe the open $\epsilon$-ball $B_\epsilon(\bar{\omega}_n)$ around a support point of $\bar{\mu} \in M(\Omega)$ and define 
\begin{equation*}
    \Omega_\epsilon := \bigcup_{n = 1}^{n = N} B_\epsilon (\bar{\omega}_n)
\end{equation*} 
$\mu \in M (\Omega_\epsilon)$ admits the form:
\begin{align*}
	&\mu = \mu_0 + \tilde{\mu} \quad\text{ with }\quad \mu_0 = \sum_{n = 1}^N c_n \delta_{\bar{\omega}_n},
    \quad 
    \tilde{\mu} = \sum_n \tilde{c}_n \delta_{\tilde{\omega}_n} + \mu_{\cont}\\
	&\text{where}\\
	& \{\tilde{\omega}_n \}_{n = 1}^\infty = \atom \mu \backslash \atom \bar{\mu}, \\
	& c_n = \mu(\{\bar{\omega}_n\}), \tilde{c}_n = \tilde{\mu}(\{\tilde{\omega}_n\}).
\end{align*}

We want to prove for all $\mu \in B_\epsilon(\bar{\mu})$, it holds $J(\bar{\mu}) \le J(\mu)$. First we observe the quadratic fidelity term, it holds
\begin{align*}
	\frac{1}{2}\|\mathcal{N}\mu - y\|_{L^2(D, \nu)}^2 &\ge \frac{1}{2} \|\mathcal{N}_{\bar{\omega}, c} - y\|^2_{L^2(D, \nu)} + \langle \mathcal{N}[\mu_0] - y, \mathcal{N} [\tilde{\mu}] \rangle_{L^2(D, \mu)} + \frac{1}{2}\|\mathcal{N}\tilde{\mu}\|^2_{L^2(D, \nu)} \\
	& \ge \frac{1}{2} \|\mathcal{N}_{\bar{\omega}, c} - y\|^2_{L^2(D, \nu)} + \langle \mathcal{N}[\mu_0 - \bar{\mu}] - y, \mathcal{N}\tilde{\mu} \rangle_{L^2(D, \mu)} + \langle \mathcal{N} \bar{\mu} - y, \mathcal{N} \tilde{\mu}\rangle_{L^2(D, \nu)}\\
	& \ge  \frac{1}{2} \|\mathcal{N}_{\bar{\omega}, c} - y\|^2_{L^2(D, \nu)} - \|\mathcal{N}\|^2 \epsilon \|\tilde{\mu}\|_{M(\Omega)} + \langle \bar{p}, \tilde{\mu}\rangle.
\end{align*}

For the penalty term, it holds
\begin{align*}
	\Phi(\mu) = \Phi(\mu_0) + \Phi(\tilde{\mu}) = \sum_{n = 1}^N \phi(|c_n|) + \sum_n \phi(|\tilde{c}_n|) + \int_\Omega d|\mu_{\cont}|.
\end{align*}

Choosing $c_n = \bar{c}_n$, the cost thus satisfies
\begin{align*}
	J(\mu) \ge J_{\bar{\omega}}(\bar{c}) - \epsilon \|\mathcal{N}\|^2 \|\tilde{\mu}\|_{M(\Omega)} + \underbrace{\alpha \Phi(\tilde{\mu}) + \langle \bar{p}, \tilde{\mu} \rangle}_{= (1)}.
\end{align*}

For $(1)$ we estimate
\begin{align*}
	\Phi(\tilde{\mu}) + \langle \bar{p}, \tilde{\mu} \rangle &= \int_\Omega \alpha + \bar{p}(\omega) d|\mu|_{\cont} + \sum_{n} \alpha \phi(|\tilde{c}_n|) + \bar{p}(\tilde{w}_n) \tilde{c}_n \\
	\bigg(|\bar{p}| \le (1 - \delta) \alpha \bigg)
	&\ge  \sum_n \alpha \phi(|\tilde{c}_n|) - (1 - \delta)|\tilde{c}_n| \\
	\bigg(\text{Assumption } \ref{pen3} (1)\bigg)
	&\ge \sum_n \alpha \delta |\tilde{c}_n| -  \frac{\gamma_1 |\tilde{c}_n|^2}{2},
\end{align*}
since $\tilde{c}_n$ is small, the $l^2$ term is dominated by the $l^1$ norm, we have 
\begin{align*}
	\Phi(\tilde{\mu}) + \langle \bar{p}, \tilde{\mu}\rangle \ge (\alpha \delta - \frac{\alpha\gamma_1}{2}) \|\tilde{c}\|_{l^1} \ge (\alpha \delta - \frac{\alpha\gamma_1}{2}) \|\tilde{\mu}\|_{M(\Omega)} .
\end{align*}
Putting the estimates together, we have
\begin{align*}
	J(\mu) &\ge J(\bar{\mu}) - \bigg(\epsilon \|\mathcal{N}\|^2 + (\alpha \delta - \frac{\alpha\gamma_1}{2}) \bigg) \|\tilde{\mu}\|_{M(\Omega)}
\end{align*} 
Choose $\epsilon$ such that  $\epsilon \|\mathcal{N}\|^2 + (\alpha \delta - \frac{\alpha\gamma_1}{2}) = 0$ we get the result. 
\end{proof}








\section{algorithms}
In this section, we introduce the optimization algorithm. In the first step, we generate the data for training and validation. In the second step, we insert neurons. And finally we trained the network with nonconvex penalty to get a sparse network. The process is repeated a fixed time. 

Denote $\mathcal{W}(D)$ as the space of functions that can be represented by measures with bounded variation, i.e., 
\begin{align*}
	&\mathcal{W}(D):= \bigg\{ f \in L^2(D, \nu) :\quad f(x) = \mathcal{N}[\mu] \text{ for some } \mu \in M(\Omega) \bigg\}, \\
	&\text{with},\\
	&\|f\|_{W(D)}: =\min_{\mu \in M(\Omega)} \|\mu\|_{M(\Omega)}\quad \text{ such that } f = \mathcal{N}[\mu].
\end{align*} 

The algorithm is divided into 3 parts:
\begin{enumerate}
	\item data generation 
	\item greedy insertion 
	\item redundant neurons elimination
\end{enumerate}
\subsection{data generation}
We come back to the original problem (\ref{opt}), the problem can be written as 
\begin{align*}
	&  \inf_{u, y} J(y, u)\\
	& \text{s.t.}\\
	& E(y, u) = 0,
\end{align*} 
with 
\begin{align*}
	E(y, u) = 
	\begin{bmatrix}
		\frac{d}{dt} y - \bigg(f(y) + g(y) u \bigg)\\
		y(0) - x 
	\end{bmatrix}.
\end{align*}
Assuming the $y$ can be written the function of $u$, define $U:= L^2(t_0, T)$, $U_{ad}:= \{u \in U:  \hat{E}(u):= E(y(u), u) = 0 \}$,  we observe the reduced problem
\begin{align*}
	& \inf_{u \in U_{ad}} \hat{J}(u)
\end{align*}
We assume the function $y(u)$ of $u$ is well-defined, the first variation at $\bar{u}$ evaluated at direction $u \in U_{ad}$ is
\begin{align*}
	&\langle \hat{J}_u(\bar{u}), u \rangle_{U^*, U} = \langle y_u^*(\bar{u}) \cdot  J_y(y(\bar{u}), \bar{u}) + J_u(y(\bar{u}), \bar{u}), u \rangle_{U^*, U},
\end{align*}
where we denote $*$ as the adjoint operator. Since $y(u)$ is implicitly defined in general and the derivative is difficult to solve, we leverage the Lagrange argument and implement the adjoint approach. Define the space $\mathcal{E}: = \ran(E)$ and $p \in  \mathcal{E}^*$. Observe
\begin{align*}
	L(p, y, u) = J(y, u) + \langle p, E(y, u) \rangle_{\mathcal{E}^*, \mathcal{E}}.
\end{align*}
Taking the partial derivative w.r.t. $y$, it holds
\begin{align*}
	&  \langle L_y(y, u), d \rangle_{Y^*, Y} = \langle J_y, d \rangle_{Y^*, Y} +   \langle E_y^*(y, u) p, d \rangle_{Y^*, Y}  \quad \text{for all } d \in Y,
\end{align*}
At the optimal pair $(y^*, u^*)$ we have $L_y(y^*, u^*) = 0$. The corresponding costate $p^*$ satisfies
\begin{equation}
\label{adjoint equation}
 	 p^*  = - E^{-*}_y \cdot  J_y \quad \text{in } \mathcal{E} 
\end{equation}
where $E_y^{-*}:= (E_y^*)^{-1}$. 
Plug $p^*$ in the Lagrangian to get the gradient w.r.t. u:
\begin{align*}
	\langle L_u(y^*, u^*, p^*), u \rangle_{U^*, U} &= \langle J_u(y^*, u^*, p^*), u \rangle_{U^*, U} + \langle E_u^*(y^*, u^*) \cdot p^*, u \rangle_{U^*, U}\\
    &= \langle 2 \beta u^* - g(y^*) p^*, u\rangle_{U^*, U} . 
\end{align*}
$p^*$ can be calculated by (\ref{adjoint equation}) and integration by parts: 
\begin{equation*}
\left\{
\begin{aligned}
	& p^*_t + \partial_y(f + g \cdot u)^T p^* = -l_y\\
	& p^*(T) = 0
\end{aligned}
\right.
\end{equation*}
For the implemented data-driven approach, it is essential to have the causality-free initial dataset $\{x^j, V(0, x^j), \nabla V(0, x^j)\}_{j = 1}^N$, where by causality-free we mean the sample are generated without approximated valuation at nearby points. \cite{kang2021algorithms} provides a comprehensive survey about generating the initial dataset. In \cite{azmi2021optimal} an adjoint approach is used to compute the optimal value and the reduced gradient (Barzilai-Borwein updates). 





\subsection{greedy insertion}
In \cite{bach2017breaking}, a conditional gradient algorithm for finitely many samples is presented. To avoid the constraint, like in \cite{pieper2022nonconvex} we implement the stereographic projection:
\begin{align*}
	\Psi: \mathbb{S}^d \rightarrow \mathbb{R}^d, \quad z \mapsto \bigg(\frac{2z}{1 + \|z\|^2}, \frac{1- \|z\|^2}{1 + \|z\|^2} \bigg).
\end{align*} 
Denote $K$ the number of the samples, the optimal neural net representation of the true target function is supported on at most $K$ atoms.  

In each iteration $f_t$, we calculate $\bar{f}_t$ such that 
\begin{align*}
	\bar{f_t} \in \arg \min_{f \in \mathcal{W}(D)} \langle f, J'(f_t) \rangle_{L^2(D, \nu)}. 
\end{align*}

 Denote $g = - J'(f_t) \in L^2(D, \nu)$ for $f_t(x) = \int_{\Omega} \varphi_\omega^t(x) d\mu_{f_t}(\omega)$. The conditional gradient algorithm relies on solving at each iteration  the "Frank-Wolf step", namely, we maximize: 
\begin{equation}
	\Omega \ni \omega \mapsto |p_t(\omega)|= \bigg|\frac{1}{K} \sum_{k = 1}^K \sigma(x_k; \omega) g_t(x_k) \bigg|. 
\end{equation}


To find the global maximal is a NP-hard problem. However heuristic approaches are available for greedy insertion.


\subsection{nonconvex penalized algorithm}
[TBD] as in Pieper.



\section{generalization bound and convergence rate}
\begin{equation}
	L(\hat{f}) - \inf_{f \in L^2(D)} L (f) \le \bigg[ \inf_{f \in L^2(D)} L(f) - \inf_{f \in \mathcal{W}(D)} L(f) \bigg] + 2 \sup_{ f \in \mathcal{W}(D)} |l(f) - L(f)| + |l(\hat{f}) - \inf_{f \in \mathcal{W}(D)} l(f)|.
\end{equation}


\section{numerical Example}
\subsection{Van der Pol  oscillator}
We consider the optimal control of the Van der Pol oscillator expressed as
\[
\min_{u \in L^2(0, T; \mathbb{R})} \int_0^T \big( y_1^2(t) + y_2^2(t) + \beta u^2(t) \big) \, dt
\]
subject to
\[
\begin{cases}
\frac{dy_1}{dt} = y_2, \\
\frac{dy_2}{dt} = -y_1 + y_2(1 - y_1^2) + u, \\
(y_1(0), y_2(0)) = (x_1, x_2),
\end{cases}
\]

The open loop problem is: 

The HJB equation at $t = 0$ is:
\begin{align*}
	\partial_t V(0, x)  -\frac{1}{4\beta} \nabla V^T(0, x) g(x)g^T(x) \nabla V(0,x) + \nabla V(0, x)^T f(x) +  l(x) = 0
\end{align*}


\section{Might be useful later }
\subsection{integral representation of convex functionals on the space of measure}
For $f: \Omega \rightarrow  (-\infty, +\infty]$ a convex function, we define the functional
\begin{align*}
	F: L^1(\Omega, \lambda) &\rightarrow \mathbb{R}\\
	 v &\mapsto \int_\Omega f(\frac{d\mu}{d\lambda}).
\end{align*}
The functional above admits a weak* l.s.c. envelope, i.e., 
\begin{align*}
	\liminf_{\mu_n \overset{*}{\rightharpoonup}\mu} F(\mu) \ge \sup_{u \in C(\Omega)} \langle \mu, u\rangle - \int_\Omega f^*(u) d\lambda.
\end{align*} 

\begin{definition}[the infimal convolution]
\label{inficonv}
	Let $f$ and $g$ be 2 closed proper convex function on $\mathbb{R}^n$, the infimal convolution of $f$ and $g$ is defined as 
	\begin{align*}
		(f \triangledown g)(v) = \inf_{x \in \mathbb{R}^n}( f(x) + g(v - x)),
	\end{align*}
	with $\dom(f \triangledown g) = \dom(f) + \dom(g)$. 
\end{definition}

\begin{proposition}
Let $f$ and $g$ be 2 closed proper convex function on $\mathbb{R}^n$, it holds
\begin{align*}
	&f \triangledown f^\infty = f, \\
	&g \triangledown g = g \quad\quad\quad \text{for $g$ subadditive and $g(0) = 0$}.
\end{align*}	
\end{proposition}

We define the indicator function  $\mathbbm{1}_H(u)$ of a closed and convex set $H \subset C(\Omega) $ with nonempty interior:
\begin{equation*}
	\mathbbm{1}_{H}(u) =  
	\left\{
	\begin{aligned}
		&0  &u \in H\\
		&+\infty &u \notin H.
	\end{aligned}
	\right.
\end{equation*} 
The Fenchel conjugate is calculated as 
\begin{align*}
	\mathbbm{1}_H^*(u^*) = \sup_{u \in H} \langle u^*, u \rangle
\end{align*}

 The following result transform the conjugation into infimal convolution. 

\begin{theorem}[infimal convolution in the functional space]
Let $\Omega$ be a compact subset of $\mathbb{R}^{d+ 1}$. Given $C(\Omega)$ with its dual $M(\Omega)$. Let $H \subset C(\Omega)$ closed and convex and denote $J$ the linear functional
\begin{align*}
	J: L^1(\Omega, m) \rightarrow \mathbb{R}, \quad u \mapsto \int_\Omega j( u(\omega)) dm(\omega). 
\end{align*}
If every $\mu \in M(\Omega)$ admits the Lebesgue-Nikodym decomposition $\mu = \frac{d\mu}{d\lambda} \cdot \lambda + \mu^s$. Then it holds for $m = |\lambda| + \mu^s$,
\begin{align*}
	[J + \mathbbm{1}_H]^*( \mu ) &= \int_\Omega [J^* \triangledown \mathbbm{1}_H^*]^{**} ( \mu ) dm = \int_\Omega j^* \nabla  \mathbbm{1}_{\Gamma}^* \bigg(\frac{d\mu}{d\lambda}(\omega)\bigg) d\mu + \int_\Omega \mathbbm{1}_{\Gamma}^*\bigg(\mu^s(\omega)\bigg)d\mu^s,
\end{align*}
with $\Gamma:= \cl \{ \ran(u) : u \in H\} \cap \dom(J)$.
\end{theorem}


\newpage
\bibliographystyle{plain}
\bibliography{ref.bib}
\end{document}




